\section{Lean 4 verification: the first two bridges}
\label{sec:lean-bridges-01-02}
\noindent\begin{minipage}[t]{0.44\textwidth}
\textbf{English: BRIDGE-001}

For any sets $A,B$ of elements of a common type, equality is equivalent
to mutual inclusion:
\[A=B\iff(A\subseteq B\land B\subseteq A).\]
If $A=B$, both inclusions follow. Conversely, two-way membership gives
equality by set extensionality.

\medskip\textbf{English: BRIDGE-002}

Set inclusion is reflexive ($A\subseteq A$), transitive
($A\subseteq B\subseteq C\Rightarrow A\subseteq C$), and
antisymmetric ($A\subseteq B\land B\subseteq A\Rightarrow A=B$).
For the final property, reuse BRIDGE-001.

\medskip\textit{Formal correspondence:} \texttt{Set $\alpha$} represents
sets whose members belong to an arbitrary type $\alpha$.
The three order laws are stated explicitly in Lean.
The code is checked separately by the pinned Lean toolchain.
\end{minipage}\hfill
\begin{minipage}[t]{0.53\textwidth}
\textbf{Lean 4 / Mathlib}

\lstinputlisting[language=LeanFour,basicstyle=\ttfamily\tiny,
breaklines=true,columns=fullflexible,keepspaces=true,frame=single]
{../../formal/lean/Con2physFormal/Bridges.lean}
\end{minipage}
